\chapter{Geltungsbereich und Grundbegriffe}

Dieser Leitfaden ergänzt Band~4 durch eine ausgewählte Übersicht
wichtiger Definitionen und Sätze über \emph{totale Relationen}.
Grundlage sind die Mengenlehre ohne Auswahlaxiom (ZF) und das
Totalitätsaxiom aus Band~4. Beweise und Beweisskizzen entfallen.
Das Format folgt den Hauptresultate-Leitfäden zu Band~3,
Band~28 und Band~37. Eine vollständige Aufzählung aller bekannten
oder ableitbaren Aussagen wird nicht beansprucht.

\section{Welche Aussagen gehören hierher?}

Jeder Satz beschreibt Totalität selbst, eine unmittelbare Folge
davon oder eine Bedingung für ihre Erhaltung. Allgemeine Gesetze
über Relationen, die mit Totalität nichts zu tun haben, werden
nicht als eigener Stoff aufgenommen. Mengenoperationen und
geordnete Paare aus Band~3 dienen als Werkzeuge.

Die behandelten Relationen tragen keine weiteren Strukturaxiome.
Insbesondere werden weder Eindeutigkeit der zugeordneten Elemente
noch Reflexivität, Symmetrie oder Transitivität vorausgesetzt.
Funktionen bleiben in den Bänden~5--8; Auswahlprinzipien in den
Bänden~9 und~16; Äquivalenzrelationen und Ordnungen in den
Bänden~11--16. Endlichkeit, Folgen, Graphen und algebraische
Strukturen werden ebenfalls nicht entwickelt.

Bedingungen innerhalb einzelner Sätze, etwa $C\subseteq A$ oder
$D\cap R_a\ne\varnothing$, sind ausdrücklich genannte
Voraussetzungen des jeweiligen Resultats. Sie sind keine
zusätzlichen Grundaxiome für alle betrachteten Relationen.
Die Aussagen gelten auch für unendliche Mengen und schließen
leere Mengen ein, soweit nichts anderes angegeben ist.
Es wird kein Auswahlaxiom vorausgesetzt oder verneint.

\section{Totalität im Sinn von Band 4}

\begin{definition}[Totale Relation]\label{b04guide:def:total}
Für Mengen $A,B,R$ bedeutet
\[
 \TR(R;A,B)\quad\Longleftrightarrow\quad
 \bigl(R\subseteq A\times B\ \land\
       \forall a\in A\,\exists b\in B\,((a,b)\in R)\bigr).
\]
Dann heißt $R$ eine \emph{totale Relation von $A$ nach $B$}.
$A$ ist ihre Ausgangsmenge und $B$ ihre Zielmenge.
Dies ist die Schreibweise $\TotRel{R,A,B}$ aus Band~4.
\end{definition}

Der Ausdruck bedeutet hier stets \emph{linkstotal}.
Für $A=B$ ist auch \emph{seriell} gebräuchlich.
Die andersartige Vergleichbarkeitsbedingung
$\forall a,b\in A\,(aRb\lor bRa)$ gehört nicht zu dieser
Definition. Die Begriffsunterscheidung findet sich auch
in \cite[Abschnitt~5.3, Fußnote~5]{meseguer}.
Die Schreibweise $aRb$ bedeutet ausschließlich $(a,b)\in R$.

\begin{definition}[Fasern und erreichte Mengen]\label{b04guide:def:fiber}
Für $R\subseteq A\times B$, $a\in A$ und $C\subseteq A$ seien
\[
 \begin{aligned}
 R_a&=\{b\in B\mid(a,b)\in R\},\\
 R[C]&=\{b\in B\mid\exists a\in C\,((a,b)\in R)\}.
 \end{aligned}
\]
$R_a$ heißt die \emph{Faser über $a$}; $R[C]$ ist die Menge
der von $C$ aus erreichten Elemente.
Insbesondere ist $R[\{a\}]=R_a$.
Die Schreibweisen bezeichnen Mengen, keine Funktionswerte.
\end{definition}

\begin{definition}[Menge aller totalen Relationen]\label{b04guide:def:all}
Für feste Mengen $A,B$ sei
\[
 \TRset(A,B)=
 \{R\in\mathcal P(A\times B)\mid \TR(R;A,B)\}.
\]
Diese Gesamtheit ist eine Menge. Die Notation
$R\in\TRset(A,B)$ und das Prädikat $\TR(R;A,B)$ sind gleichbedeutend.
\end{definition}

\section{Orientierung}

Kapitel~\ref{b04guide:ch:criteria} charakterisiert Totalität und
klärt die leeren Fälle. Kapitel~\ref{b04guide:ch:sets} beschreibt,
wann Mengenoperationen Totalität erhalten.
Kapitel~\ref{b04guide:ch:combine} behandelt das Verketten und
Koppeln totaler Relationen. Kapitel~\ref{b04guide:ch:models}
gibt Modelle aus Mengenfamilien.

Die Nummern mit Präfix $4H$ gehören ausschließlich zu diesem
Leitfaden. Die Nummerierung des Fachbands bleibt unabhängig davon.
